\chapter[When the model itself is incomplete]{When the model itself is incomplete: missingness, misspecification, and research priorities}
\label{sim:model-risk}

Preserving all histories inside a model does not protect against a missing mechanism outside it. A feasible family can be exact and still omit a guarantee, a customer concentration or a financing restriction. This chapter examines how to represent that limitation without abandoning useful analysis. The aim is to make conditional conclusions precise and to direct further work toward omissions that can change decisions.

\section{Unknown value versus absent mechanism}
An unknown interest rate can be represented by a variable if the contract form is known. An unknown existence of a guarantee requires an alternative contract structure. A mechanism the language cannot express is a different limitation. These cases should not share one generic null.

A missing value within a supported form can often be bounded or inferred. A missing structural alternative can change the dimension and event logic of the family. An unsupported mechanism may require a new model component before any bound is meaningful. The result should identify which kind of gap it faces.

For example, allowing a loan rate between zero and twenty per cent does not cover a missed conversion right. A wider rate interval cannot reproduce the ownership change caused by conversion. Numerical sensitivity over existing parameters is not a substitute for structural coverage.

\section{The disclosure package defines an evidence boundary}
A company filing can incorporate other documents by reference or discuss risks without quantifying them. A reconstruction based only on primary statements may omit material terms in notes, exhibits or subsequent events. The evidence package should state which documents and passages were considered.

Coverage has stages: retrieval, parsing, interpretation and model admission. A page can be retrieved but not parsed; parsed but not understood; understood but outside the current executable profile. These distinctions matter for the completeness claim.

An unresolved passage should remain linked to relevant model areas. A guarantee discussion may block a claim that all contingent obligations are represented. It need not block a simple historical revenue reconciliation. Scope should be query-sensitive rather than all-or-nothing.

\section{Missingness can be informative}
A missing disclosure may result from immateriality, reporting exemptions, extraction failure or deliberate omission within a strategic disclosure model. These causes have different implications. Treating every missing value as zero assumes the most benign explanation. Treating every missing value as arbitrarily large can make analysis useless.

The first task is to identify the source of missingness where possible. If a table was not extracted from a scanned page, the remedy is data processing. If the filing does not disclose a contract term, the uncertainty is substantive. If a disclosure is legally required under certain facts, inference from absence requires actual authority and evidence about compliance; it cannot be assumed.

A probabilistic model of missingness can be added when justified. It should state whether missingness depends on hidden exposure or actor behavior. An assumption of missing at random is not a universal default for corporate disclosure. Without support, retain alternative explanations rather than fabricate a likelihood.

\section{Conflicting sources}
An annual report, interim filing and debt exhibit can disagree. The disagreement can reflect different dates, scopes, currencies, amendments or errors. A model that takes the latest number without reconciling context can discard valid historical information.

Source precedence should depend on the claim. A later amendment may supersede future contract terms while leaving past payments unchanged. A restated financial statement may revise a comparative amount. A management presentation may use a nonstandard metric rather than contradict a filing.

The evidence view should preserve both assertions and the interpretation linking them. If conflict remains unresolved, the family may contain alternative source-admission views or return an explicit contradiction. It should not average incompatible facts merely because both are numeric.

\section{Measurement error and reporting discretion}
A discrepancy between predicted and reported amounts can arise from rounding, estimation, extraction error or an incorrect economic history. A single Gaussian residual model collapses these mechanisms. It may be convenient statistically but can imply implausible support and independence.

Rounding often gives bounded intervals. Extraction errors can be discrete shifts in sign or scale. Accounting estimates can have structured dependence on assumptions and information. Reporting discretion may be strategic. The observation model should reflect the mechanism relevant to the claim.

A robust likelihood or contamination model can be considered, but its parameters need provenance and sensitivity analysis. It should not become a way to make every contradiction disappear. A large residual may be evidence that the model omitted a transaction rather than merely noisy reporting.

\section{Structural alternatives}
Suppose revenue is modeled as price times quantity. A different model includes subscription renewals, usage charges and cancellation options. Both may fit historical totals. A demand shock can produce different future responses because the contract structures differ.

The family can include alternative structural models with shared observations and explicit scope. Comparing their conditional outcomes identifies model risk. Assigning posterior weights requires model priors and likelihoods; otherwise the alternatives remain unweighted.

A structural alternative should be concrete enough to simulate or bound. A vague statement that anything could happen is not useful. The modeler should identify plausible omitted mechanisms, their source support and the smallest discriminating evidence. This makes uncertainty actionable without claiming completeness.

\section{Behavioral parameter uncertainty}
A collection-delay parameter estimated from historical data can be uncertain and unstable. A recession can change customer behavior, credit terms or the company's enforcement policy. Treating the estimate as a fixed fact understates uncertainty.

Separate estimation uncertainty within a behavioral form from uncertainty about the form itself. A confidence interval around a constant collection rate does not cover a regime switch in which major customers stop paying. Conversely, a broad regime set may be unnecessarily pessimistic for a short, well-understood contract schedule.

The scenario should state whether behavior is held fixed, estimated conditionally, changed by a stress rule or optimized. Each treatment answers a different question. An empirical policy estimated under one financing regime may not transfer after covenants bind.

\section{Extrapolation}
A model calibrated in low-rate periods can be applied to high-rate stress. Contractual cash formulas may remain valid if terms are known, while behavioral relationships may not. The extrapolation status should be component-specific.

A floating coupon formula with a known cap can be evaluated outside historical rates because the contract supplies the rule. A regression of sales on rates may have no support in that region. A neural surrogate trained only below the cap may also fail even though the underlying contract formula is simple.

The model should distinguish extrapolation of an economic law, a behavioral estimate and a numerical approximation. Each requires different evidence. A global warning that the scenario is extreme is less useful than identifying which components are unsupported and how that affects the result.

\section{Unknown unknowns and explicit coverage}
A model cannot enumerate mechanisms it has not conceived. It can, however, maintain a coverage map organized by assets, obligations, rights, entities, timing and information. Review can ask whether material disclosures have a place in that map.

Stress testing the representation itself is useful. Introduce an unexpected guarantee, a delayed tax payment or a conversion clause and see whether the system can preserve it as unresolved rather than dropping it. A robust data model should fail visibly when a new mechanism is outside its profile.

This is different from claiming that a checklist proves completeness. The coverage map supports systematic discovery. It remains extensible as new cases arise. The manuscript's failure catalogue should be read in that spirit.

\section{Sensitivity to the family boundary}
Suppose a result is safe for every admitted history. The conclusion can still depend on the family's bounds. If floating share is limited to 0.4 by a weak assumption rather than a source, expanding it to one can change safety.

A useful analysis identifies active assumptions: those whose relaxation changes the conclusion or its bound. This can be done through alternative runs, dual information in suitable optimization problems or counterexample search. The method should not infer economic importance solely from a numerical coefficient.

An assumption may be inactive for one query and decisive for another. A maturity bound irrelevant to next month's interest can dominate a one-year refinancing stress. Sensitivity records should therefore be tied to query and horizon.

\section{Partial identification of a decision}
The evidence may not identify every parameter but still identify the preferred action. If one policy dominates another across all compatible histories, further historical detail is unnecessary for that comparison. Conversely, a narrow parameter interval can still straddle a decision threshold.

This distinction directs research effort. The goal is not always to recover the true journal uniquely. It may be enough to establish a robust decision or a bounded exposure. A history family is useful precisely because it can show which unresolved details matter.

The system should therefore report both parameter uncertainty and decision sensitivity. A large uncertainty set need not imply indecision, and a small set need not imply safety. The geometry relative to policies and thresholds is what matters.

\section{Research priorities from counterexamples}
A verified adverse witness can reveal the missing fact that drives risk. If the witness requires a near-term foreign-currency maturity, the next source target is a currency-maturity cross-classification. If it requires a guarantee, inspect the guarantee terms. If it requires refinancing refusal, investigate committed capacity and draw conditions.

The witness should be accompanied by competing interpretations and checked consequences. This is useful internal preparation before asking a professional a specific question. An unresolved label alone is not a review-ready brief.

The repository's internal-preparation requirement therefore fits the method: finish feasible source retrieval, calculations and competing interpretations first. No professional engagement or expenditure is authorized by the existence of an open accounting or legal question. Ordinary reader feedback is a separate need.

\section{Model revision without losing history}
When a missing mechanism is added, old results may change. Preserve the old model and result with their scope, then create a revised version. This allows users to distinguish a forecast error under the old model from a correction to the target definition.

A new guarantee can expand the family or add obligations. A new disclosure can narrow it. A corrected accounting mapping can change reported outcomes without changing economic cash. These revisions have different meanings and should not all be called data refreshes.

Dependency tracking should identify which results need recomputation. A source affecting one loan should not require rebuilding unrelated evidence, but its effects may propagate through liquidity, default and group guarantees. The dependency graph must follow economic consequences as well as direct formulas.

\section{Communicating model risk}
A useful explanation states the result, the assumptions that drive it and the specific unresolved mechanisms that could change it. It does not bury the conclusion under generic warnings. For example, ``cash remains above 14 for all retained histories provided the facility remains drawable after customer default'' identifies a meaningful condition.

If that condition is not established, say so directly and show the alternative result. Avoid presenting a conditional safe outcome as unconditional safety. Equally, do not describe every unresolved mechanism as inevitable disaster. The family should distinguish possible, necessary and unassessed consequences.

The reader should be able to see whether more evidence, a different policy or a better computation would resolve the issue. These are different next steps. A numerical approximation gap cannot be closed by another disclosure if the model already knows the relevant fact; a missing contract clause cannot be closed by more solver iterations.

\section{A defensible stopping point}
A query can be complete within a declared scope even when the entire company model is incomplete. Completion means that the requested target has been computed with the stated guarantees and that material excluded mechanisms are disclosed. It does not mean all uncertainty has disappeared.

A broader product-readiness claim requires broader coverage and evidence. The 150--300-page treatment is a documentation deliverable, not proof of that readiness. The manuscript remains a candidate whose assumptions, sources and proposed methods require continued implementation and human review.

This distinction permits useful progress. It avoids both false certainty and paralysis. The system can answer bounded questions accurately, retain uncertainty honestly and identify the concrete work needed to support stronger conclusions.
